\documentclass[a4paper,10pt]{article}

\usepackage[margin=0.5in]{geometry}
\usepackage{enumitem}
\usepackage{titlesec}
\usepackage{hyperref}
\usepackage{fontawesome5}

%----------FORMATTING----------
\titleformat{\section}
{\large\bfseries}
{}
{0em}
{}
[\titlerule]

\titlespacing{\section}{0pt}{5pt}{3pt}

\setlist[itemize]{
    leftmargin=16pt,
    itemsep=0pt,
    topsep=1pt,
    parsep=0pt
}

%----------HYPERLINKS----------
\hypersetup{
    colorlinks=true,
    urlcolor=blue,
    linkcolor=black,
    pdfborder={0 0 0}
}

\begin{document}

%----------HEADER----------
\begin{center}

{\Huge\textbf{Deepak Gupta}}\\[5pt]

\small
\faPhone\ +91-9458777101
\quad
\href{mailto:mrdeepak.g11@gmail.com}{\faEnvelope\ mrdeepak.g11@gmail.com}
\quad
\href{https://linkedin.com/in/imdeepx11}{\faLinkedin\ linkedin.com/in/imdeepx11}
\quad
\href{https://github.com/imdeepx11}{\faGithub\ github.com/imdeepx11}

\end{center}

%----------SUMMARY----------
\section{Summary}

Electronics and Communication Engineering student with hands-on experience in embedded systems, Arduino-based hardware projects, and full-stack AI application development. Skilled in integrating software and hardware for real-world solutions, with exposure to industrial railway signalling systems. Seeking opportunities to apply embedded systems and software development skills in a challenging engineering role.

%----------SKILLS----------
\section{Skills}

\begin{itemize}
    \item \textbf{Programming:} Python, C, Java (Basic), JavaScript (Basic)
    \item \textbf{Frameworks \& Tools:} FastAPI, React, Arduino IDE, Git, Vercel, Render
    \item \textbf{Embedded Systems:} Embedded C, Arduino (UNO/Nano), Sensor Integration, PWM Control, Bluetooth Communication
    \item \textbf{Core Electronics:} Digital Electronics, Basic Electronics, Circuit Design
    \item \textbf{Other:} REST API Development, SQLite, AI-assisted development and productivity tools
\end{itemize}

%----------PROJECTS----------
\section{Projects}

\textbf{NEXORA AI -- Intelligent Document Processing \& Workflow Automation}
\hfill
\href{https://github.com/imdeepx11/NEXORA-AI}{\textit{[GitHub]}}
\enspace
\href{https://nexora-ai-imdeepx11.vercel.app/}{\textit{[Live Demo]}}

\begin{itemize}
    \item Conceptualized and developed a full-stack AI web application for intelligent document processing and workflow automation, leveraging AI-assisted development tools throughout the build.
    \item Designed and integrated an end-to-end pipeline: document upload (PDF/DOCX/TXT) $\rightarrow$ AI classification $\rightarrow$ structured field extraction $\rightarrow$ risk scoring (0--100) $\rightarrow$ automated workflow routing with audit trail.
    \item Orchestrated the tech stack -- \textbf{Python}/\textbf{FastAPI} backend, \textbf{React 19}/\textbf{Tailwind CSS}/\textbf{Vite} frontend -- and managed cloud deployment on \textbf{Vercel} and \textbf{Render}.
    \item Gained practical understanding of \textbf{REST API} design, \textbf{SQLite} data persistence, \textbf{PyMuPDF} document extraction, and modern full-stack architecture through hands-on implementation.
\end{itemize}

\vspace{2pt}

\textbf{Sentinel Rover: Radar-Assisted Bluetooth-Controlled Autonomous Vehicle}
\hfill
\href{https://drive.google.com/drive/folders/1aF45gX7f-_E48fL2EvL6yxTDeOj20PwY?usp=sharing}{\textit{[Project Videos]}}

\begin{itemize}
    \item Developed an Arduino-based rover featuring Bluetooth-controlled navigation and radar-assisted real-time object detection.
    \item Implemented semi-autonomous target tracking using sensor feedback loops and embedded control logic.
    \item Integrated motor driver module, PWM-based speed control, buzzer feedback system, and wireless command processing.
\end{itemize}

\vspace{2pt}

\textbf{Sun Tracking System Using Arduino}

\begin{itemize}
    \item Built an automatic solar tracking system using Arduino UNO, LDR sensors, and a servo motor to dynamically orient a solar panel toward maximum sunlight exposure.
    \item Implemented real-time dual-axis light intensity comparison logic to improve simulated panel efficiency.
\end{itemize}

\vspace{2pt}

\textbf{Automatic Gate Opener Using Arduino}

\begin{itemize}
    \item Designed an automated gate control system using Arduino UNO, proximity sensor, and servo motor with object detection logic for automatic open/close actuation.
\end{itemize}

%----------INTERNSHIP----------
\section{Internship}

\textbf{Signal \& Telecom Training Centre, N.E. Railway, Gorakhpur}
\hfill Jun. 2026 -- Jul. 2026

\textit{Summer Intern -- Railway Signalling \& Telecommunications Systems}

\begin{itemize}
    \item Completed a one-month industrial internship focused on railway signalling and telecommunication infrastructure.
    \item Gained practical exposure to signalling equipment, inter-station communication systems, and railway safety mechanisms.
\end{itemize}

%----------EDUCATION----------
\section{Education}

\textbf{JSS Academy of Technical Education, Noida} \hfill 2023 -- 2027 (Expected)

B.Tech in Electronics and Communication Engineering \hfill CGPA: 6.97 / 10

\vspace{3pt}

Class XII -- ICSE Board: \textbf{80.2\%} (2022) \hfill Class X -- ICSE Board: \textbf{89.16\%} (2020)

%----------CERTIFICATIONS----------
\section{Certifications}

\begin{itemize}
    \item VLSI for Beginners -- NIELIT Calicut (MeitY, Govt.\ of India)
    \item Python Programming -- Udemy
\end{itemize}

%----------ACHIEVEMENTS----------
\section{Achievements \& Leadership}

\begin{itemize}
    \item Appointed as Technical Squad Member of \textbf{Quanta} -- Electronics Department Technical Society, JSS Noida
    \item Served as \textbf{Co-Core Team Member} of Quizzoc Society; managed event planning, logistics, and team coordination
    \item Successfully completed government-certified VLSI training under \textbf{NIELIT Calicut} (MeitY, Govt.\ of India)
\end{itemize}

\end{document}
